% Part: proof-theory
% Chapter: proof-search
% Section: introduction

\documentclass[../../../include/open-logic-section]{subfiles}

\begin{document}

\olfileid{pt}{ps}{int}

\olsection{పరిచయం}

\Log{G3c}, \Log{N1i} వంటి \tetoken{నిరూపణ}{!!{proof}} పద్ధతుల ఆక్సియమ్‌లు, నిగమన నియమాలు ఏ సీక్వెంట్ లేదా \tetoken{సూత్ర}{!!{formula}s} వృక్షాలు సరైన \tetoken{నిరూపణలుగా}{!!{proof}s} లెక్కించబడతాయో నిర్వచిస్తాయి. సీక్వెంట్‌లు లేదా \tetoken{సూత్రాల}{!!{formula}s} వృక్షం ఇచ్చినప్పుడు (ప్రతి దశలో వాడిన నియమాలు, ఊహల గుర్తులు, వాటిని విడుదల చేసే నియమాల గుర్తులు వంటి అదనపు సమాచారం కూడా ఉండవచ్చు), అది సరైన \tetoken{నిరూపణా}{!!{proof}} కాదా అని ఆక్సియమ్‌లు, నియమాల నిర్వచనాలతో తనిఖీ చేయవచ్చు. ఇచ్చిన సీక్వెంట్‌కు లేదా ఇచ్చిన ఊహల నుంచి ఇచ్చిన \tetoken{సూత్రానికి}{!!{formula}} \tetoken{ఒక నిరూపణను}{!!a{proof}} \emph{కనుగొనడం} వేరొక సమస్య; సాధారణంగా అది మరింత కష్టం. ఇదే \emph{నిరూపణ అన్వేషణ} సమస్య.

నిరూపణ అన్వేషణకు చాలా సరళమైన పద్ధతి ఉంది. \tetoken{సూత్రాలను}{!!{formula}s} పరిమిత వర్ణమాలలోని చిహ్నాల వరుసలుగా భావించవచ్చు. సీక్వెంట్‌లు, సీక్వెంట్ లేదా \tetoken{సూత్ర}{!!{formula}s} వృక్షాలను కూడా \tetoken{సూత్రాల}{!!{formula}s} వరుసలుగా భావించవచ్చు (వృక్ష శాఖలను సూచించడానికి ప్రత్యేక కుండలీకరణలను వాడవచ్చు). చిహ్నాల వరుస సరైన \tetoken{నిరూపణను}{!!{proof}} సూచిస్తుందో లేదో ఆక్సియమ్‌లు, నియమాలతో యాంత్రికంగా తనిఖీ చేయవచ్చు. అందువల్ల \emph{సాధ్యమైన అన్ని} సరైన \tetoken{నిరూపణలను}{!!{proof}s} యాంత్రికంగా ఉత్పత్తి చేయవచ్చు: \tetoken{నిరూపణలను}{!!{proof}s} సూచించే వర్ణమాల నుంచి అన్ని చిహ్న వరుసలను ఉత్పత్తి చేసి, ప్రతి అభ్యర్థి సరైన \tetoken{నిరూపణను}{!!{proof}} సూచిస్తుందో తనిఖీ చేయాలి. సరళమైన అన్వేషణ పద్ధతి ఇది: ఇచ్చిన సీక్వెంట్‌కు (లేదా ఇచ్చిన \tetoken{సూత్రానికి}{!!{formula}}, ఇచ్చిన ఊహల్లోనే తెరిచి ఉన్న ఊహలు ఉండేలా) \tetoken{నిరూపణ}{!!{proof}} దొరికే వరకు అన్ని సరైన \tetoken{నిరూపణలను}{!!{proof}s} ఉత్పత్తి చేస్తూ పోవాలి.

చేతితో \tetoken{ఒక నిరూపణను}{!!a{proof}} కనుగొనడానికి వాడే సాధారణ పద్ధతి మరింత ఉపయోగకరం: చివరి సీక్వెంట్‌తో మొదలుపెట్టి \tetoken{ఒక సూత్రాన్ని}{!!a{formula}} ఎంచుకుని, నిగమన నియమాన్ని "వెనుకకు" వర్తింపజేసి ఒకటి లేదా ఎక్కువ పూర్వపక్షాలను ఉత్పత్తి చేయాలి. వాటికి \tetoken{నిరూపణలు}{!!{proof}s} ఉంటే, చివరి అనుమితితో కలిపి ఇచ్చిన సీక్వెంట్‌కు \tetoken{ఒక నిరూపణను}{!!a{proof}} పొందవచ్చు. సహజ నిగమన పద్ధతిలో \tetoken{ఒక నిరూపణను}{!!a{proof}} కనుగొనడానికీ ఇలాంటి పద్ధతినే వాడతాం, అయితే అది కొంత క్లిష్టం.

కానీ ఈ పద్ధతి తప్పకుండా విజయాన్ని ఇవ్వదు: తప్పు నియమాన్ని ఎంచుకున్నా, \tetoken{సూత్రాలను}{!!{formula}s} తప్పు క్రమంలో ఎంచుకున్నా, ముందుకు వెళ్లలేని చోటికి చేరవచ్చు. ఇది పూర్తిగా నిర్దేశించబడిన పద్ధతీ కాదు: ప్రతి దశలో ఎంచుకోవడానికి ఒకటికంటే ఎక్కువ \tetoken{సూత్రాలు}{!!{formula}} లేదా వెనుకకు వర్తింపజేయడానికి ఒకటికంటే ఎక్కువ నియమాలు ఉండవచ్చు. \emph{నిరూపణ అన్వేషణ అల్గోరిథం} అంటే పూర్తిగా నిర్దేశించబడిన పద్ధతి; \tetoken{ఒక నిరూపణ}{!!a{proof}} ఉంటే దాన్ని తప్పకుండా కనుగొంటుంది.

కొన్ని \tetoken{నిరూపణ}{!!{proof}} పద్ధతుల్లో అన్వేషణ అల్గోరిథం ఇతర పద్ధతులకంటే సరళంగా ఉంటుంది. సీక్వెంట్ పద్ధతుల్లో \tetoken{ఒక సూత్రం}{!!a{formula}} యొక్క \tetoken{ప్రధాన సంయోజకం}{!!{main operator}}, అది ఎడమవైపునా కుడివైపునా ఉందనే విషయం, ఏ నియమాన్ని వెనుకకు వర్తింపజేయాలో నిర్ణయిస్తాయి. ఉదాహరణకు, $!A \land
!B$ కుడివైపున ఉన్న $\Gamma \Sequent \Delta, !A \land !B$కు \tetoken{ఒక నిరూపణను}{!!a{proof}} కనుగొనాలంటే, \RightR{\land}ను వెనుకకు వర్తింపజేసి $\Gamma \Sequent \Delta, !A$, $\Gamma \Sequent \Delta, !B$ అనే రెండు పూర్వపక్షాలను ఉత్పత్తి చేయాలి. కానీ పద్ధతి \Log{G1c} అయితే సంకోచనం వల్ల మరిన్ని అవకాశాలు వచ్చి సమస్య క్లిష్టమవుతుంది: \RightR{\Contraction}ను వెనుకకు వర్తింపజేసి $\Gamma \Sequent \Delta, !A \land
!B, !A \land !B$ అనే పూర్వపక్షాన్ని కూడా ఉత్పత్తి చేయవచ్చు. తరువాత $\Gamma \Sequent
\Delta, !A \land !B, !A \land !B, !A \land !B$, ఇలా కొనసాగించవచ్చు. \Log{G2c}లో మరింత క్లిష్టం. \RightR{\land}ను వెనుకకు వర్తింపజేయాలంటే $\Gamma
= \Gamma_1 \cup \Gamma_2$ అయ్యేలా $\Gamma_1$, $\Gamma_2$లను, $\Delta
= \Delta_1 \cup \Delta_2$ అయ్యేలా $\Delta_1$, $\Delta_2$లను ఊహించి ఎంచుకుని, $\Gamma_1 \Sequent \Delta_1, !A$, $\Gamma_2 \Sequent \Delta_2, !B$ పూర్వపక్షాలతో ప్రయత్నించాలి. తప్పుగా ఎంచుకుంటే తరువాత ముందుకు వెళ్లలేని స్థితి రావచ్చు; అప్పుడు మొదటికి వెనక్కి వచ్చి వేరే $\Gamma_1$, $\Gamma_2$, $\Delta_1$, $\Delta_2$లను ఎంచుకోవాలి. \tetoken{సూత్రం}{!!{formula}} $!A \lor !B$ అయితే \RightR{\lor}కు ఉన్న రెండు అవకాశాలలో ఒకదాన్ని ఎంచుకోవాలి; అది $\Gamma \Sequent \Delta,
!A$ లేదా $\Gamma \Sequent \Delta, !B$ అనే పూర్వపక్షాన్ని ఇస్తుంది. తప్పు ఎంపిక వల్ల మళ్లీ వెనక్కి రావాల్సి వస్తుంది.

\Log{G3c}కు ఈ సమస్యలు లేవు. దానిలో నిర్మాణ నియమాలు లేవు; నిగమన నియమ ఎంపికను \tetoken{ప్రధాన సంయోజకం}{!!{main operator}}, \tetoken{సూత్రం}{!!{formula}} ఉన్న వైపు నిర్ణయిస్తాయి. కాబట్టి \Log{G1c}, \Log{G2c}లకంటే \Log{G3c} నిరూపణ అన్వేషణకు అనుకూలం. విజయవంతమైన అన్వేషణ \Log{G3c}లో \tetoken{ఒక నిరూపణను}{!!a{proof}} ఇస్తుంది; ఆ \tetoken{నిరూపణను}{!!{proof}} \Log{G1c}, \Log{G2c} లేదా \Log{N1c}లోకి అనువదించవచ్చు. అందువల్ల \Log{G3c}కు అన్వేషణ అల్గోరిథం, \Log{G3c} \tetoken{నిరూపణలను}{!!{proof}s} ఇతర పద్ధతుల \tetoken{నిరూపణలుగా}{!!{proof}s} మార్చే అనువాద అల్గోరిథం ఉంటే, ఆ పద్ధతుల నిరూపణ అన్వేషణ సమస్య కూడా పరిష్కరించబడుతుంది.

అన్వేషణ అల్గోరిథం తప్పకుండా సరైన నిర్ణయం ఇస్తుందని ఎలా తెలుసుకోగలం? అది \tetoken{ఒక నిరూపణను}{!!a{proof}} కనుగొనడంలో విఫలమైతే, \tetoken{నిరూపణ}{!!{proof}} అసలు ఉండదని చూపాలి. చెడు అల్గోరిథాన్ని ఎంచుకుంటే, నిరూపణ ఉన్నప్పటికీ అది కనుగొనలేని సందర్భాలు రావచ్చు. సరళమైన అల్గోరిథానికి ఈ సమస్య లేదు; అది సాధ్యమైన అన్ని \tetoken{నిరూపణలనూ}{!!{proof}s} వెతుకుతుంది. అయితే నిరూపణ అన్వేషణ అల్గోరిథం సమర్థంగా, కనీస శ్రమతో పనిచేయాలని కోరుకుంటాం.

నిరూపణ అన్వేషణ అల్గోరిథం సరైన నిర్ణయం ఇస్తుందని చూపడానికి కీలక భావం ఇది: అల్గోరిథం విఫలమైతే, అది విఫలమయ్యే విధానం నుంచే ఆ సీక్వెంట్‌కు \tetoken{ఒక నిరూపణ}{!!a{proof}} ఉండదని చూపాలి. దానికోసం ఆ సీక్వెంట్ సర్వసత్యం కాదని చూపుతాం. సంప్రదాయ ప్రథమ క్రమ తర్కంలో $\Gamma \Sequent \Delta$ సర్వసత్య సీక్వెంట్ కావాలంటే, ప్రతి \tetoken{నిర్మాణంలో}{!!{structure}} $\Gamma$లో కనీసం ఒక \tetoken{సూత్రం}{!!{formula}} అసత్యం కావాలి లేదా~$\Delta$లో కనీసం ఒక \tetoken{సూత్రం}{!!{formula}} సత్యం కావాలి. \Log{G3c} \emph{సబబైనది}; అంటే అది సర్వసత్య సీక్వెంట్‌లను మాత్రమే \tetoken{నిరూపిస్తుంది}{!!{prove}s}. \tetoken{నిరూపణలపై}{!!{proof}s} ఆగమనంతో దీన్ని స్థాపిస్తాం: ఆక్సియమ్‌లు స్పష్టంగా సర్వసత్యమైనవి; నియమ పూర్వపక్షాలు సర్వసత్యమైతే నిష్కర్ష కూడా సర్వసత్యమే. అల్గోరిథం సరైన నిర్ణయం ఇస్తుందని చూపడానికి, $\Gamma \Sequent \Delta$కు \tetoken{ఒక నిరూపణ}{!!a{proof}} కోసం అన్వేషణ విఫలమైతే, $\Gamma$లోని అన్ని \tetoken{సూత్రాలు}{!!{formula}s} సత్యం,~$\Delta$లోని అన్ని \tetoken{సూత్రాలు}{!!{formula}s} అసత్యం అయ్యే \tetoken{ఒక నిర్మాణాన్ని}{!!a{structure}} నిర్వచించగలమని చూపుతాం. దీనివల్ల \Log{G3c} \emph{సంపూర్ణమైనదని} కూడా స్థాపించబడుతుంది: $\Gamma \Sequent \Delta$ సర్వసత్యమైతే దానికి \Log{G3c}-\tetoken{నిరూపణ}{!!{proof}} ఉంటుంది. విఫలమయ్యే ప్రతి అన్వేషణ $\Gamma \Sequent \Delta$ సర్వసత్యం కాదని చూపుతుంది కాబట్టి, సర్వసత్య సీక్వెంట్ కోసం ప్రతి అన్వేషణ విజయవంతం కావాలి.

\end{document}
